\subsection{\texorpdfstring{\(\wedge^{p}(\mathbf{M})\) 与 \(\wedge^{n - p}(\mathbf{M}^{*})\) 之间的同构（\(n\) 维自由模的情形）}{\textbackslash wedge\^{}\{p\}(\textbackslash mathbf\{M\}) 与 \textbackslash wedge\^{}\{n - p\}(\textbackslash mathbf\{M\}\^{}\{*\}) 之间的同构（n 维自由模 M）}}

命题 12. 设 \(\mathbf{M}\) 是一个维数为 \(n\) 的自由 A-模；令 \(e \in \bigwedge^{n}(\mathbf{M})\) 是构成 \(\bigwedge^{n}(\mathbf{M})\) 的一组基的元素，并设 \(e^{*}\) 是 \(\bigwedge^{n}(\mathbf{M}^{*})\) 中使得 \(\{(-1)^{n(n-1)/2} \theta_{\bigwedge}(e^{*})\}\) 是 \(\{e\}\) 在 \((\bigwedge^{n}(\mathbf{M}))^{*}\) 中的对偶基的元素。设 \(\phi: \bigwedge(\mathbf{M}^{*}) \to \bigwedge(\mathbf{M})\) 是映射 \(z \mapsto z \lrcorner e^{*}\) ，而 \(\phi': \bigwedge(\mathbf{M}^{*}) \to \bigwedge(\mathbf{M})\) 是映射 \(z^{*} \mapsto z^{*} \lrcorner e\) 。设 \(\phi_{p}\) （相应地 \(\phi_{p}'\) ）是 \(\phi\) （相应地 \(\phi'\) ）在 \(\bigwedge^{p}(\mathbf{M})\) （相应地 \(\bigwedge^{p}(\mathbf{M}^{*})\) ）上的限制。那么：

\begin{enumerate}
\def\labelenumi{(\roman{enumi})}
\item
  映射 \(\phi\) 是一个左 \(\bigwedge (\mathbf{M})\) -模同构，而映射 \(\phi'\) 是一个左 \(\bigwedge (\mathbf{M}^*)\) -模同构；此外，映射 \(\phi\) 和 \(\phi'\) 互为逆映射。
\item
  映射 \(\phi_p\) 是 A-模 \(\bigwedge^p (\mathbf{M})\) 到 A-模 \(\bigwedge^{n - p}(\mathbf{M}^*)\) 上的同构，而映射 \(\phi_p'\) 是 A-模 \(\bigwedge^p (\mathbf{M}^*)\) 到 A-模 \(\bigwedge^{n - p}(\mathbf{M})\) 上的同构。
\item
  若记 \(\mathbf{B}(u,v^{*}) = \langle u,\theta_{\wedge}(v^{*})\rangle\) （对 \(u\in \bigwedge (\mathbf{M})\) 和 \(v^{*}\in \bigwedge (\mathbf{M}^{*})\) ），则对 \(u^{*}\in \bigwedge^{p}(\mathbf{M}^{*})\) 和 \(v^{*}\in \bigwedge^{n - p}(\mathbf{M}^{*})\)
\end{enumerate}

\[
\mathbf {B} \left(\phi_ {p} ^ {\prime} (u ^ {*}), v ^ {*}\right) = (- 1) ^ {p (n - p)} \mathbf {B} \left(u ^ {*}, \phi_ {n - p} ^ {\prime} (v ^ {*})\right).\tag{77}
\]

\(\phi\) 是 \(\bigwedge (\mathbf{M})\) -线性的以及 \(\phi'\) 是 \(\bigwedge (\mathbf{M}^*)\) -线性的这一事实，由公式 \((u \wedge v) \lrcorner e^* = u \lrcorner (v \lrcorner e^*)\) 和 \((u^* \wedge v^*) \lrcorner e = u^* \lrcorner (v^* \lrcorner e)\) 得出（第 6 号，公式 (37)，利用 \(\theta_{\wedge}\) 是 \(\bigwedge (\mathbf{M}^*)\) 到 \(\bigwedge (\mathbf{M}^*)\) 的反代数上的同构这一事实）。另一方面，存在一个基 \((e_i)_{1 \leq i \leq n}\) （ \(\mathbf{M}\) 的基），使得

\[
e = e _ {1} \wedge e _ {2} \wedge \dots \wedge e _ {n} \quad \text { and } \quad e ^ {*} = (- 1) ^ {n (n - 1) / 2} e _ {1} ^ {*} \wedge e _ {2} ^ {*} \wedge \dots \wedge e _ {n} ^ {*},
\]

（其中 \((e_i^*)\) 是 \((e_i)\) 的对偶基。我们记 \(\mathbf{I} = [1, n]\) ；由 (76) 可知，对每个子集 \(\mathbf{J}\) （属于 \(\mathbf{I}\) ，含 \(p\) 个元素），

\[
\left\{ \begin{array}{l} \phi (e _ {J}) = (- 1) ^ {n (n - 1) / 2 + p (p - 1) / 2} \rho_ {J, I - J} e _ {I - J} ^ {*} \\ \phi^ {\prime} (e _ {J} ^ {*}) = (- 1) ^ {p (p - 1) / 2} \rho_ {J, I - J} e _ {I - J}. \end{array} \right.\tag{78}
\]

这证明了 \(\phi\) 和 \(\phi'\) 是双射；此外， \(\rho_{J,I-J}\rho_{I-J,J}=(-1)^{p(n-p)}\) （ \(\S 7\) ，第 8 号，公式 (21)）；由于数

\[
\frac {n (n - 1)}{2} + \frac {p (p - 1)}{2} + \frac {(n - p) (n - p - 1)}{2} + p (n - p) = n (n - 1)
\]

是偶数，因此得出 \(\phi\) 和 \(\phi'\) 互为逆映射。最后，为证明 (77)，只需取 \(u^{*} = e_{J}^{*}\) 和 \(v^{*} = e_{I - J}^{*}\) ；验证同样由 \(\theta_{\wedge}\) 的定义、公式 (78) 及关系式 \(\rho_{J,I - J}\rho_{I - J,J} = (-1)^{p(n - p)}\) （§ 7，第 8 号，公式 (21)）得出。注意，对 \(u^{*}\in \bigwedge^{p}(\mathbf{M}^{*})\) 和 \(v^{*}\in \bigwedge^{n - p}(\mathbf{M}^{*})\) , \(\mathbf{B}(\phi_{p}^{\prime}(u^{*}), v^{*})\) 是 \(u^{*} \wedge v^{*}\) 关于基 \(\{e^{*}\}\) （ \(\bigwedge^{n}(\mathbf{M}^{*})\) 的基）的系数，至多差一个符号。

命题 13. 在命题 11 的假设和记号下，对 A-模 M 的每个自同态 g，

\[
(\det g) \phi = \bigwedge (^ {t} g) \circ \phi \circ \bigwedge (g).\tag{79}
\]

显然 \(\bigwedge(^{t}g) = \theta_{\wedge}^{-1}\circ (^{t}\bigwedge (g))\circ \theta_{\wedge}\) ；因为 \(\bigwedge (g)\) 是代数 \(\bigwedge (\mathbf{M})\) 的自同态，并且根据定义，对所有

\[
z \in \bigwedge (\mathbf {M}), \quad \theta_ {\wedge} (\bigwedge (g) (z) \lrcorner e ^ {*}) = \theta_ {\wedge} (e ^ {*}) \llcorner \theta_ {\wedge} (\bigwedge (g) (z)),
\]

我们由第 6 号的公式 (42) 推出

\[
\begin{array}{r l} ((\theta_ {\wedge} ^ {- 1} \circ (^ {t} \bigwedge (g)) \circ \theta_ {\wedge}) \circ \phi \circ \bigwedge (g)) (z) & = \theta_ {\wedge} ^ {- 1} (^ {t} \bigwedge (g) (\theta_ {\wedge} (e ^ {*})) \llcorner z) \\ & = z \lrcorner (\bigwedge (t g) (e ^ {*})) = (\det g) (z \lrcorner e ^ {*}) = (\det g) \phi (z) \end{array}
\]

这里用到了 § 8，第 4 号，命题 8。

推论。对 E 的每个自同构 g，

\[
\bigwedge(^ {t} g ^ {- 1}) = (\det g) ^ {- 1} \phi \circ (\bigwedge (g)) \circ \phi^ {- 1}.\tag{80}
\]

